\begin{frame}{우리 팀 발표를 리뷰하기 — 발표 전 셀프-리뷰}
  \begin{itemize}
    \item ML1 8.2의 결론을 그대로: \kb{네가 심사할 때 던질 수 있는
      가장 날카로운 질문을, 네 발표 \emph{전에} 스스로 던져보고 답을
      준비해두는 것}이 최선의 전략.
    \item §5가지 패턴의 5개 질문을 \textbf{자기 팀의 보고서}에 적용:
      \begin{enumerate}
        \item 시드 3개의 이동평균 표가 \textbf{보고서에} 있는가
        \item 평가 프로토콜 문장(탐험 0? 에피소드 수? 모든 시드에 동일)이
          있는가
        \item 무작위 정책의 리턴이 \textbf{실측 숫자}로 적혀 있는가
        \item $\lambda\cdot$보상스케일 민감성을 \emph{한 줄이라도}
          검증했는가
        \item 불안정한 시드·구간이 ``잘 안 된 부분'' 절에 \emph{사라지지
          않고} 있는가
      \end{enumerate}
    \item 이 5가지에 답이 없다면, \textbf{발표 전에} 해결해두는 것이
      발표 중에 3점 리뷰를 받는 것보다 \textbf{항상 낫다}.
  \end{itemize}
\end{frame}
